\documentclass[UTF8, a4paper]{ctexbook}
\usepackage{geometry}
\usepackage{setspace}
\usepackage{booktabs}
\usepackage{longtable}
\usepackage{listings}
\usepackage{xcolor}
\usepackage{hyperref}
\usepackage{caption}
\usepackage{fancyhdr}
\usepackage{enumitem}
\usepackage{titlesec}
\usepackage{amsmath}
\usepackage{graphicx}

% 页面设置
\geometry{left=2.5cm, right=2.5cm, top=2.5cm, bottom=2.5cm}

% 页眉高度修正
\setlength{\headheight}{13pt}

% 页眉页脚
\pagestyle{fancy}
\fancyhf{}
\fancyhead[C]{\thepage}
\renewcommand{\headrulewidth}{0pt}

% 段落间距
\setlength{\parskip}{0.5\baselineskip}
\setlength{\parindent}{2em}

% 标题格式
\titleformat{\section}{\Large\bfseries}{\thesection}{1em}{}
\titleformat{\subsection}{\large\bfseries}{\thesubsection}{1em}{}
\titleformat{\subsubsection}{\normalsize\bfseries}{\thesubsubsection}{1em}{}

% ============ 定义 YAML 和 Properties 语言（简化，避免特殊字符） ========================
\lstdefinelanguage{yaml}{
	keywords={true,false,null,y,n},
	keywordstyle=\color{blue}\bfseries,
	basicstyle=\ttfamily\small,
	sensitive=false,
	comment=[l]{\#},
	string=[b]",
	string=[b]',
}

\lstdefinelanguage{properties}{
	keywords={},
	keywordstyle=\color{blue},
	basicstyle=\ttfamily\small,
	comment=[l]{\#},
	string=[b]",
	string=[b]',
}

% 全局代码风格
\lstset{
	basicstyle=\ttfamily\small,
	keywordstyle=\color{blue},
	commentstyle=\color{gray},
	stringstyle=\color{red},
	showstringspaces=false,
	frame=single,
	breaklines=true,
	backgroundcolor=\color{gray!5},
	tabsize=2,
	numbers=left,
	numberstyle=\tiny\color{gray},
	stepnumber=1,
	numbersep=5pt,
}

% 超链接颜色
\hypersetup{
	colorlinks=true,
	linkcolor=blue,
	urlcolor=blue,
	citecolor=blue,
}

% 文档信息
\title{Apache DolphinScheduler}
\author{运维技术团队}
\date{\today}

\begin{document}
	
\maketitle
\thispagestyle{empty}
\newpage
	
\tableofcontents
\thispagestyle{empty}
\newpage
	
\setcounter{page}{1}
	
\chapter{前言}
\addcontentsline{toc}{section}{前言}
	
本文档为 Apache DolphinScheduler 在生产环境下的部署与运维提供全面指导，主要面向非集群部署（单机或轻量级多进程）场景，核心诉求是\textbf{数据持久化}，即确保调度元数据、任务状态、执行日志等关键信息不因服务重启而丢失。文档涵盖了三种主流部署方案（Standalone + 外部数据库、Docker Compose、Kubernetes Helm），并从配置调优、监控告警、安全加固、备份恢复、故障排查及版本升级等方面给出详细建议。希望本文档能成为您在生产环境中落地 DolphinScheduler 的实用参考。
	
\chapter{概述}
	
\section{文档目的}
	
Apache DolphinScheduler 是一个分布式、可扩展的\textbf{可视化 DAG 工作流任务调度平台}。在高可用生产环境中，通常推荐部署多 Master/Worker 集群，但在资源受限或业务初期小规模场景下，用户往往选择单机或轻量级方式，同时要求数据持久化。本文档聚焦于这种“非集群 + 持久化”的部署模式，提供经过验证的完整实施方案。
	
\section{核心组件架构}
	
DolphinScheduler 采用 Master-Worker 去中心化设计，各组件职责如下：
	
\begin{longtable}{p{3cm}p{10cm}}
	\toprule
	\textbf{组件} & \textbf{职责} \\
	\midrule
	MasterServer & 解析 DAG、监控任务状态与容错、生成工作流实例、派发任务到 Worker \\
	WorkerServer & 具体执行任务，向注册中心（ZooKeeper）注册自身 \\
	ApiServer  & 提供 RESTful API 服务，供前端 UI 调用 \\
	AlertServer & 告警服务，支持邮件、钉钉、企业微信等插件 \\
	Registry (ZooKeeper) & Master 选举、Worker 注册与发现、分布式锁 \\
	\bottomrule
\end{longtable}
	
\section{硬件配置基线}
	
DolphinScheduler 支持 64 位 Intel x86-64 硬件平台。单机部署建议不低于 4 核 CPU、8GB 内存。各节点推荐配置如下：
	
\begin{longtable}{p{3cm}p{3cm}p{3cm}p{3cm}}
	\toprule
	\textbf{节点类型} & \textbf{CPU} & \textbf{内存} & \textbf{磁盘} \\
	\midrule
	Master & 4核+ & 8GB+ & 50GB+ \\
	Worker & 8核+ & 16GB+ & 50GB+ \\
	API   & 4核+ & 8GB+  & 50GB+ \\
	Alert & 2核+ & 4GB+  & 50GB+ \\
	\bottomrule
\end{longtable}
	
\chapter{方案选型对比}
	
\begin{longtable}{p{3.5cm}p{2.5cm}p{2.5cm}p{3.5cm}}
	\toprule
	\textbf{特性} & \textbf{Standalone + 外置DB} & \textbf{Docker Compose} & \textbf{Kubernetes Helm} \\
	\midrule
	部署复杂度 & 低   & 中   & 高   \\
	资源占用   & 低（单进程） & 中（多容器） & 中高 \\
	稳定性     & 一般 & 较好 & 好   \\
	数据持久化 & 外部数据库 & 外部DB+数据卷 & 外部DB+PV \\
	高可用     & 不支持 & 不支持 & 不支持（单点） \\
	适用场景   & 开发/测试、<20工作流 & 小规模生产 & 已有K8s环境 \\
	维护成本   & 低   & 中   & 高   \\
	\bottomrule
\end{longtable}
	
\textbf{推荐建议}：对于大多数“非集群但需持久化”的生产场景，\textbf{Docker Compose 方案} 是性能、稳定性和维护成本的最佳平衡点。
	
\chapter{Standalone + 外部数据库(方案一)}
	
\section{前置准备}
	
\begin{itemize}
	\item JDK 1.8+ 已安装并配置环境变量。
	\item MySQL 5.7+ 或 PostgreSQL 10+ 已安装并启动。
	\item 创建专用数据库 \texttt{dolphinscheduler}，字符集 \texttt{utf8mb4}。
	\item 创建专用数据库用户并授予该库全部权限。
\end{itemize}
	
\section{安装包与驱动}
	
\begin{lstlisting}[language=bash, caption=下载并解压二进制包]
	# 下载二进制发行包（以 3.2.2 为例）
	wget https://downloads.apache.org/dolphinscheduler/3.2.2/apache-dolphinscheduler-3.2.2-bin.tar.gz
	
	# 解压
	tar -zxvf apache-dolphinscheduler-3.2.2-bin.tar.gz -C /opt/
	cd /opt/apache-dolphinscheduler-3.2.2-bin
	
	# 下载 MySQL 驱动并复制到 standalone-server 目录
	wget https://repo1.maven.org/maven2/mysql/mysql-connector-java/8.0.33/mysql-connector-java-8.0.33.jar
	cp mysql-connector-java-8.0.33.jar standalone-server/libs/standalone-server/
\end{lstlisting}
	
\section{修改数据库配置}
	
编辑 \texttt{standalone-server/conf/application.yaml}，将数据源改为外部数据库。
	
\textbf{MySQL 配置示例}：
	
\begin{lstlisting}[language=yaml]
	spring:
	datasource:
	driver-class-name: com.mysql.cj.jdbc.Driver
	url: jdbc:mysql://127.0.0.1:3306/dolphinscheduler?useUnicode=true&characterEncoding=UTF-8&serverTimezone=Asia/Shanghai&useSSL=false
	username: ds_user
	password: your_secure_password
	hikari:
	connection-test-query: SELECT 1
	minimum-idle: 5
	maximum-pool-size: 50
	pool-name: dsHikariCP
	auto-commit: true
	idle-timeout: 30000
	max-lifetime: 1800000
\end{lstlisting}
	
\textbf{PostgreSQL 配置示例}：
	
\begin{lstlisting}[language=yaml]
	spring:
	datasource:
	driver-class-name: org.postgresql.Driver
	url: jdbc:postgresql://127.0.0.1:5432/dolphinscheduler
	username: ds_user
	password: your_secure_password
	hikari:
	connection-test-query: SELECT 1
	minimum-idle: 5
	maximum-pool-size: 50
\end{lstlisting}
	
\section{初始化数据库}
	
\begin{lstlisting}[language=bash]
	sh tools/bin/upgrade-schema.sh
\end{lstlisting}
	
\section{启动与验证}
	
\begin{lstlisting}[language=bash]
	# 启动 Standalone 服务
	./bin/dolphinscheduler-daemon.sh start standalone-server
	
	# 查看日志确认启动成功
	tail -f logs/dolphinscheduler-standalone-server.log
	
	# 访问 Web UI: http://<your_ip>:12345
	# 默认用户名/密码: admin/dolphinscheduler123
\end{lstlisting}
	
\section{局限性}
	
\begin{itemize}
	\item Standalone 模式内置单机 ZooKeeper，存在单点故障风险。
	\item 官方不建议用于超过 20 个工作流的场景。
	\item 日志和任务输出文件默认保存在本地，需自行备份。
	\item 默认 H2 内存数据库在重启后数据丢失，故\textbf{必须更换为外部数据库}。
\end{itemize}
	
\chapter{Docker Compose 部署(方案二)}
	
\section{环境要求}
	
\begin{itemize}
	\item Docker 20.10+ 和 Docker Compose 1.28+。
	\item 主机至少 4GB 内存（推荐 8GB）。
\end{itemize}
	
\section{获取部署文件}
	
\begin{lstlisting}[language=bash]
	git clone https://github.com/apache/dolphinscheduler.git
	cd dolphinscheduler/deploy/docker
\end{lstlisting}
	
\section{部署步骤}
	
\begin{lstlisting}[language=bash]
	# 初始化数据库（首次或升级时执行）
	docker compose --profile schema up -d
	
	# 启动所有服务
	docker compose --profile all up -d
	
	# 查看运行状态
	docker compose ps
	
	# 查看日志
	docker compose logs -f
\end{lstlisting}
	
\section{数据持久化说明}
	
Compose 文件默认定义了三个 Docker 卷，确保数据在容器删除后依然保留：
	
\begin{longtable}{p{5cm}p{8cm}}
	\toprule
	\textbf{卷名} & \textbf{用途} \\
	\midrule
	\texttt{postgresql-data}    & PostgreSQL 数据目录 \\
	\texttt{zookeeper-data}     & ZooKeeper 数据目录 \\
	\texttt{dolphinscheduler-log} & 任务执行日志目录 \\
	\bottomrule
\end{longtable}
	
\section{自定义配置}
	
通过环境变量调整 JVM 参数和线程数，例如在 \texttt{docker-compose.yaml} 中：
	
\begin{lstlisting}[language=yaml]
	environment:
	- MASTER_HEAP_MEMORY=2g
	- WORKER_HEAP_MEMORY=4g
	- WORKER_EXEC_THREADS=50
\end{lstlisting}
	
\section{与其他服务合并}
	
若需将 DolphinScheduler 与其他应用（如 Nginx、Redis）统一管理，可采用多 Compose 文件合并：
	
\begin{lstlisting}[language=bash]
	# 方式一：命令行指定多个文件
	docker compose -f docker-compose.yaml -f my-other-services.yaml up -d
	
	# 方式二：使用 COMPOSE_FILE 环境变量
	export COMPOSE_FILE=docker-compose.yaml:my-other-services.yaml
	docker compose up -d
	
	# 方式三：重命名为 docker-compose.override.yml（自动加载）
	mv my-other-services.yaml docker-compose.override.yml
	docker compose up -d
\end{lstlisting}
	
\chapter{Kubernetes Helm 部署(方案三)}
	
\section{前置条件}
	
\begin{itemize}
	\item Kubernetes 1.18+ 集群，已配置 StorageClass。
	\item Helm 3.0+ 客户端。
	\item 外部数据库（MySQL/PostgreSQL）。
\end{itemize}
	
\section{部署步骤}
	
\begin{lstlisting}[language=bash]
	# 添加 Helm 仓库
	helm repo add dolphinscheduler https://apache.github.io/dolphinscheduler
	helm repo update
	
	# 创建 values-prod.yaml（见下方配置示例）
	# 部署
	helm install dolphinscheduler dolphinscheduler/dolphinscheduler -f values-prod.yaml -n dolphinscheduler --create-namespace
	
	# 验证状态
	kubectl get pods -n dolphinscheduler
\end{lstlisting}
	
\section{生产环境 values-prod.yaml 配置}
	
\begin{lstlisting}[language=yaml]
	# ==================== 镜像配置 ====================
	# 建议将镜像提前推送到私有仓库
	image:
	registry: your-private-registry.com
	tag: "3.2.2"
	
	# ==================== 数据库配置 ====================
	# 生产环境务必关闭内置 PostgreSQL
	postgresql:
	enabled: false
	
	externalDatabase:
	type: mysql  # 或 postgresql
	host: "your-db-host"
	port: 3306
	username: "ds_user"
	password: "secure_password"
	database: "dolphinscheduler"
	
	# ==================== 注册中心配置 ====================
	zookeeper:
	enabled: false
	
	externalZookeeper:
	servers: "zk1:2181,zk2:2181,zk3:2181"
	
	# ==================== 副本数配置（非集群）====================
	api:
	replicaCount: 1
	master:
	replicaCount: 1
	worker:
	replicaCount: 1
	alert:
	replicaCount: 1
	
	# ==================== 持久化配置 ====================
	worker:
	logPersistence:
	enabled: true
	storageClass: "standard"
	size: 10Gi
	
	# ==================== LDAP 认证（可选）====================
	authentication:
	ldap:
	enabled: true
	url: "ldap://ldap.example.com:389"
	baseDn: "dc=example,dc=com"
	userDn: "uid={0},ou=people,dc=example,dc=com"
	
	# ==================== 共享存储（可选）====================
	persistence:
	sharedStorage:
	enabled: true
	path: "/opt/soft"
	existingClaim: "nfs-shared-storage"
\end{lstlisting}
	
\section{踩坑经验（基于奇虎360实践）}
	
\begin{itemize}
	\item MySQL 驱动必须软链到每一个模块：\texttt{dolphinscheduler-tools}、\texttt{dolphinscheduler-master}、\texttt{dolphinscheduler-worker}、\texttt{dolphinscheduler-api} 和 \texttt{dolphinscheduler-alert-server}。
	\item 基础镜像过大，建议将公共软件层拆成多阶段构建以缓存。
	\item 自编译 Jar 包未覆盖旧包时，确保构建时清理旧文件。
\end{itemize}
	
\chapter{统一配置调优指南}
	
\section{Master Server 配置}
	
位置：\texttt{master-server/conf/application.yaml}
	
\begin{longtable}{p{5cm}p{3cm}p{5cm}p{3cm}}
	\toprule
	\textbf{参数} & \textbf{默认值} & \textbf{说明} & \textbf{调优建议} \\
	\midrule
	\texttt{master.exec-threads} & 100 & 工作线程数，限制并行流程实例数量 & CPU核数×2~4 \\
	\texttt{master.pre-exec-threads} & 10 & 准备执行任务数量 & 流程多时可适当增大 \\
	\texttt{master.dispatch-task-number} & 3 & 每个批次派发任务数量 & 可调至10–20 \\
	\texttt{master.server-load-protection.enabled} & true & 是否开启系统保护策略 & 生产建议开启 \\
	\texttt{master.server-load-protection.max-system-cpu-usage-percentage-thresholds} & 0.7 & CPU阈值，超过则停止调度 & 根据实际调整 \\
	\texttt{master.worker-load-balancer-configuration-properties.type} & \texttt{DYNAMIC\_WEIGHTED\_ROUND\_ROBIN} & 负载均衡策略 & 生产推荐此策略 \\
	\bottomrule
\end{longtable}
	
\section{Worker Server 配置}
	
位置：\texttt{worker-server/conf/application.yaml}
	
\begin{longtable}{p{5cm}p{3cm}p{5cm}p{3cm}}
	\toprule
	\textbf{参数} & \textbf{默认值} & \textbf{说明} & \textbf{调优建议} \\
	\midrule
	\texttt{worker.exec.threads} & 100 & 任务执行线程数 & 根据资源调整 \\
	\texttt{worker.task.max.memory} & 4096 & 单任务最大内存(MB) & 根据任务需求调整 \\
	\texttt{worker.tenant-auto-create} & true & 是否自动创建租户 & 生产建议关闭 \\
	\texttt{worker.host-weight} & 100 & Worker 主机权重 & 多Worker差异化 \\
	\texttt{worker.worker.task.resource.limit} & false & 是否开启任务资源限制 & 建议开启 \\
	\bottomrule
\end{longtable}
	
\section{JVM 参数调优}
	
在 \texttt{bin/dolphinscheduler-daemon.sh} 中配置：
	
\begin{lstlisting}[language=bash]
	# Master 节点
	export JAVA_OPTS="-Xms4g -Xmx4g -XX:+UseG1GC -XX:MaxGCPauseMillis=20"
	
	# Worker 节点
	export JAVA_OPTS="-Xms8g -Xmx8g -XX:+UseG1GC -XX:MaxGCPauseMillis=20"
\end{lstlisting}
	
\textbf{\textcolor{red}{重要警告}}：\textbf{不要}设置 \texttt{-XX:DisableExplicitGC}，否则可能导致 Netty 内存泄漏。
	
\section{API Server 流量控制}
	
\begin{longtable}{p{5cm}p{3cm}p{5cm}}
	\toprule
	\textbf{参数} & \textbf{默认值} & \textbf{说明} \\
	\midrule
	\texttt{api.traffic.control.default-tenant-qps-rate} & 10 & 默认租户最大请求数/秒 \\
	\texttt{api.traffic.control.customize-tenant-qps-rate} & - & 自定义租户QPS限制 \\
	\bottomrule
\end{longtable}
	
\chapter{监控体系搭建}
	
\section{监控架构}
	
DolphinScheduler 使用 Micrometer 作为指标采集框架，通过 Spring Boot Actuator 端点暴露 Prometheus 格式指标：
	
\begin{itemize}
	\item 所有服务组件（Master、Worker、API、Alert）均可暴露指标。
	\item 端点地址：\texttt{http://host:port/actuator/prometheus}
	\item 默认端口：Master 5679、Worker 1235、Alert 50053、API 12345
\end{itemize}
	
\section{启用监控}
	
在各服务 \texttt{application.yaml} 中配置：
	
\begin{lstlisting}[language=yaml]
	metrics:
	enabled: true
	
	management:
	endpoints:
	web:
	exposure:
	include: prometheus,health
\end{lstlisting}
	
\section{Grafana 面板}
	
官方提供了开箱即用的 Grafana 面板配置，位于 \texttt{dolphinscheduler-meter/resources/grafana}：
	
\begin{lstlisting}[language=bash]
	cd dolphinscheduler-meter/src/main/resources/grafana-demo
	docker compose up -d
	# 访问 http://localhost:3001
\end{lstlisting}
	
\section{核心监控指标}
	
\begin{longtable}{p{5cm}p{5cm}p{3cm}}
	\toprule
	\textbf{指标名称} & \textbf{描述} & \textbf{告警阈值} \\
	\midrule
	\texttt{master.jvm.memory.used} & Master JVM 内存使用量 & >85\% \\
	\texttt{worker.thread.pool.active} & Worker 活跃线程数 & >90\% 线程池 \\
	\texttt{task.success.rate} & 任务成功率 & <95\% \\
	\texttt{database.connection.usage} & 数据库连接使用率 & >85\% \\
	\bottomrule
\end{longtable}
	
\chapter{安全加固方案}
	
\section{敏感信息泄露防护（CVE-2025-62188）}
	
Apache DolphinScheduler 3.1.x 版本存在未授权访问 Actuator 端点导致泄露数据库凭证的漏洞。修复方案如下：
	
\begin{lstlisting}[language=bash]
	# 方案一：升级到 3.2.0 或更高版本（推荐）
	# 方案二：临时限制端点暴露
	export MANAGEMENT_ENDPOINTS_WEB_EXPOSURE_INCLUDE=health,metrics,prometheus
\end{lstlisting}
	
或在 \texttt{application.yaml} 中配置：
	
\begin{lstlisting}[language=yaml]
	management:
	endpoints:
	web:
	exposure:
	include: health,metrics,prometheus
\end{lstlisting}
	
\section{配置加密}
	
支持使用 AES 对称加密对敏感属性进行加密：
	
\begin{enumerate}
	\item 生成 AES 密钥（配置 \texttt{encrypt.secret.key}）。
	\item 将明文敏感值替换为 \texttt{ENC(加密后的密文)} 格式。
	\item 系统启动时自动解密。
\end{enumerate}
	
\begin{lstlisting}[language=yaml]
	spring:
	datasource:
	password: ENC(encrypted_password_here)
\end{lstlisting}

对于更高安全要求，可集成外部 KMS。
	
\section{权限管理}
	
\begin{itemize}
	\item 严格控制管理员数量，仅系统搭建者可拥有管理员权限。
	\item 生产环境建议关闭 \texttt{worker.tenant-auto-create}，改为手动管理租户。
	\item 通过 Worker 分组（\texttt{conf/worker.properties}）实现资源隔离。
	\item 利用安全中心的租户、用户和角色实现细粒度权限控制。
\end{itemize}
	
\section{系统安全基线}
	
\begin{itemize}
	\item 关闭不必要的防火墙规则（按需开放端口）。
	\item 在 \texttt{/etc/security/limits.conf} 中为特定用户设置资源限制。
	\item 启用自动安全更新。
\end{itemize}
	
\chapter{数据备份与恢复}
	
\section{数据库备份}
	
\textbf{MySQL 备份}：
	
\begin{lstlisting}[language=bash]
	# 全量备份
	mysqldump -u root -p dolphinscheduler > /backup/dolphinscheduler_$(date +%Y%m%d).sql
	
	# 定时备份（crontab）
	0 2 * * * mysqldump -u root -p'password' dolphinscheduler > /backup/ds_$(date +\%Y\%m\%d).sql
\end{lstlisting}
	
\textbf{PostgreSQL 备份}：
	
\begin{lstlisting}[language=bash]
	pg_dump -U ds_user -h localhost dolphinscheduler > /backup/dolphinscheduler_$(date +%Y%m%d).sql
\end{lstlisting}
	
升级之前务必备份数据，防止操作错误。
	
\section{资源文件备份}
	
\begin{lstlisting}[language=bash]
	tar -czf /backup/ds_resources_$(date +%Y%m%d).tar.gz /path/to/dolphinscheduler/resources/
\end{lstlisting}
	
\section{数据恢复}
	
\begin{lstlisting}[language=bash]
	# MySQL 恢复
	mysql -u root -p dolphinscheduler < /backup/dolphinscheduler_20260101.sql
	
	# PostgreSQL 恢复
	psql -U ds_user -d dolphinscheduler < /backup/dolphinscheduler_20260101.sql
\end{lstlisting}
	
\section{资源迁移（3.2.0+）}
	
从 3.2.0 版本起资源中心重构，可指定迁移目标租户，运行一次性资源迁移脚本，所有资源将迁移到目标租户的 \texttt{.migrate} 目录。
	
\chapter{常见故障排查}
	
\section{服务启动失败}
	
\textbf{症状}：服务无法启动或启动后立即退出。
	
\textbf{排查步骤}：
	
\begin{enumerate}
	\item \textbf{端口冲突}：检查端口占用并修改。
	\begin{lstlisting}[language=bash]
		netstat -tlnp | grep 5678   # Master 端口
		netstat -tlnp | grep 12345  # API 端口
	\end{lstlisting}
	\item \textbf{内存不足}：调整 JVM 参数。
	\item \textbf{依赖服务未启动}：确保 ZooKeeper 和数据库正常运行。
	\item \textbf{JDK 版本不匹配}：确认使用 JDK 1.8+。
	\item \textbf{数据库驱动缺失}：检查驱动 JAR 是否放入正确目录。
\end{enumerate}
	
\section{任务调度异常}
	
\textbf{症状}：任务卡在“提交中”，工作流实例无法生成。
	
\textbf{排查步骤}：
	
\begin{enumerate}
	\item 检查 Master 服务状态：\texttt{jps | grep MasterServer}
	\item 查看 ZooKeeper 注册状态：\texttt{echo stat | nc localhost 2181}
	\item 检查数据库连接数和慢查询：
	\begin{lstlisting}[language=sql]
		SHOW PROCESSLIST;
		SHOW VARIABLES LIKE 'max_connections';
		EXPLAIN SELECT * FROM t_ds_process_instance WHERE state = 1;
	\end{lstlisting}
	\item 排查 MySQL 死锁（检查相关表）。
\end{enumerate}
	
\section{UI 无法登录}
	
\textbf{排查步骤}：
	
\begin{enumerate}
	\item 检查 API 服务连通性。
	\item 验证 Session 配置：
	\begin{lstlisting}[language=properties]
		server.servlet.session.timeout=3600
		server.context-path=/dolphinscheduler
	\end{lstlisting}
	\item 确认静态资源文件存在：\texttt{ls -la /path/to/ui/static/}
\end{enumerate}
	
\section{网络与 IP 地址问题}
	
\textbf{解决方案}：指定 IP 获取策略和网卡。
	
\begin{lstlisting}[language=properties]
	dolphin.scheduler.network.priority.strategy=default
	dolphin.scheduler.network.interface.preferred=eth0
\end{lstlisting}
	
\textbf{连通性测试}：
	
\begin{lstlisting}[language=bash]
	ping worker-node-ip
	telnet worker-node-ip 12345
	iptables -L -n   # 检查防火墙
\end{lstlisting}
	
\chapter{维护与升级}
	
\section{日常维护清单}
	
\begin{longtable}{p{3cm}p{10cm}}
	\toprule
	\textbf{频率} & \textbf{操作} \\
	\midrule
	每日 & 检查服务进程、查看错误日志、监控磁盘使用率 \\
	每周 & 清理过期日志、检查数据库连接池、审查任务执行情况 \\
	每月 & 数据库备份验证、性能分析、检查安全更新 \\
	每季度 & 评估版本升级、容量规划、灾难恢复演练 \\
	\bottomrule
\end{longtable}
	
\section{服务启停}
	
\begin{lstlisting}[language=bash]
	# 启动所有服务
	./bin/start-all.sh
	
	# 停止所有服务
	./bin/stop-all.sh
	
	# 单独启停某个服务
	./bin/dolphinscheduler-daemon.sh start master-server
	./bin/dolphinscheduler-daemon.sh stop worker-server
\end{lstlisting}
	
\section{升级流程}
	
\begin{enumerate}
	\item \textbf{备份数据}：数据库和资源文件。
	\item \textbf{阅读升级文档}：关注 Breaking Changes。
	\item \textbf{测试环境验证}：先升级测试环境。
	\item \textbf{执行升级脚本}：\texttt{tools/bin/upgrade-schema.sh}
	\item \textbf{验证功能}：确认所有服务正常。
	\item \textbf{生产环境升级}：选择低峰期执行。
\end{enumerate}
	
\section{版本选择建议}
	
\begin{longtable}{p{3cm}p{3cm}p{6cm}}
	\toprule
	\textbf{版本} & \textbf{推荐程度} & \textbf{说明} \\
	\midrule
	3.1.x & 不推荐 & 存在 CVE-2025-62188 安全漏洞 \\
	3.2.x & 推荐   & 修复安全漏洞，功能稳定 \\
	3.3.x & 最新   & 更多新特性，建议充分测试后使用 \\
	\bottomrule
\end{longtable}
	
\chapter{附录}
\addcontentsline{toc}{section}{附录}
	
\section*{A. 常用命令速查}

\begin{longtable}{p{5cm}p{7cm}}
	\toprule
	\textbf{操作} & \textbf{命令} \\
	\midrule
	查看服务状态 & \texttt{jps | grep -E "MasterServer|WorkerServer|ApiServer"} \\
	查看 ZooKeeper 状态 & \texttt{echo stat | nc localhost 2181} \\
	查看数据库连接 & \texttt{SHOW PROCESSLIST;} \\
	查看端口占用 & \texttt{netstat -tlnp | grep <port>} \\
	查看实时日志 & \texttt{tail -f logs/dolphinscheduler-*.log} \\
	\bottomrule
\end{longtable}
	
\section*{B. 关键配置文件位置}

\begin{longtable}{p{5cm}p{7cm}}
	\toprule
	\textbf{配置文件} & \textbf{路径} \\
	\midrule
	环境变量 & \texttt{bin/env/dolphinscheduler\_env.sh} \\
	数据库配置 & \texttt{conf/application.yaml} \\
	注册中心配置 & \texttt{conf/application.yaml} \\
	日志配置 & \texttt{conf/logback.xml} \\
	Common 配置 & \texttt{conf/common.properties} \\
	\bottomrule
\end{longtable}
	
\section*{C. 参考链接}

\begin{itemize}
	\item 官网：\url{https://dolphinscheduler.apache.org}
	\item GitHub：\url{https://github.com/apache/dolphinscheduler}
	\item 开发者邮件列表：dev@dolphinscheduler@apache.org
\end{itemize}
	
\end{document}